
I aim to quantify the degree to which eternal inflation is a generic feature of scalar inflation models.  In other words, what is the likelihood of finding ourselves in an eternal multiverse, assuming no {\it a priori} knowledge of the governing potential function or initial state?

\begin{enumerate}

\item
Elements of this analysis have been performed by many researchers at the level of individual models, but it has never been extended to systematically address the entire class.  My first task is to catalogue those published models for which inflation is known to be non-eternal for most empirically consistent sets of parameters and initial conditions.  The analytical criteria for eternal inflation are well documented; this step consists in reviewing those criteria evaluated for particular models, and diagnosing their failure in the generic instance.  Based on patterns discerned from those negative cases, I will establish robust criteria on which to base my classification of any given scalar inflation model as generically eternal or non-eternal.  

\item
Since the meaning of generic does not extend unambiguously to describe the whole class of models, the next step is to develop a sensible measure – a procedure for compartmentalizing the space of candidate models for the purpose of assigning probabilities.  The measure establishes the statistical framework in which I will assess the likelihood of eternal inflation, yet known laws do not appear to single out a preferred measure.  I must make a choice, weighing the successes of existing measures 3 with factors such as simplicity, alignment with a consistent set of physical principles, and utility in sorting the features that cosmologists are most interested in studying.  

\item
The full space of candidate models is too large to be explored analytically, so in order to make broad statistical statements about it, I must sample the space of inflation models numerically.  I will simulate the dynamics of inflation with randomly generated potentials and initial conditions, recording which models (delimited by the chosen measure) are generically eternal according to my established criteria.  Finally, I will analyze those numerical results to make a quantitative assessment of the genericity of eternal inflation.  
